Post-quantum cryptography (PQC), also known as quantum-resistant 
cryptography, aims at defining new public-key cryptographic schemes based 
on mathematical problems proven to be computationally hard against 
adversaries equipped with both classical and quantum computers. 
Recent advancements on quantum computer technology have witnessed a steep 
increase in scientific and engineering activities to deliver new solutions 
to some of the most pressing scientific and societal challenges~\cite{mckinsey23}. 
However, from a cybersecurity standpoint, the availability of a quantum computer 
with a sufficiently large number of qubits, would leave most of the 
current protocols used to guarantee the confidentiality and authenticity of communications and data storage vulnerable to cryptanalytic techniques based on Shor~\cite{DBLP:conf/focs/Shor94}-like quantum algorithms that subvert the 
security guarantees of the underlying public key cryptosystems.\\
\indent Since $2017$, several institutional and governmental bodies promoted the 
analysis of quantum-secure systems, but most stakeholders track the efforts 
of the PQC standardization contest promoted by the U.S. National Institute 
of Standard (NIST).  
Through two rounds of this evaluation process, NIST selected seven finalists
and eight alternate candidates, which were then thinned down to two standard
candidates and four alternates for a final round of cryptanalysis and possible 
standardization. 
The current standardization-elected candidate Key Encapsulation Method (KEM) is 
Kyber~\cite{Kyber}, a lattice-based cryptosystem, while NIST selected four 
candidates relying on different computationally hard problems, in order to obtain
resilience against potentially disruptive attacks on lattice-based cryptosystems.
Among the four candidates, a polynomial-time cryptanalytic attack eliminated
SIKE, leaving only three KEMs which are based based on the hardness of the 
syndrome decoding problem, i.e., Classic McEliece~\cite{ClassicMcEliece}, 
BIKE~\cite{BIKE} and HQC~\cite{HQC}.
Among these candidates, only HQC can exhibit a security proof 
based on a polynomial reduction of the solution of a computationally
hard problem to the breaking of the cryptosystem.
In particular, any successful break of HQC implies a faster way to solve the
syndrome decoding problem on a random quasi-cyclic code of rate $\frac{1}{2}$.
By contrast, Classic McEliece relies on an additional assumption, namely the indistinguishability of a hidden Goppa code from a random code, and BIKE relies 
on the indistinguishability of a hidden QC-MDPC code from a random quasi 
cyclic code.
While both arguments are currently unchallenged, there is no proof that they
are indeed corresponding to computationally hard problems.\\
\indent Hamming Quasi-Cyclic (HQC)~\cite{HQC,DBLP:journals/tit/Aguilar-Melchor18} is
a code-based IND-CCA secure KEM whose security relies on the syndrome 
decoding problem for quasi cyclic codes. It is obtained by applying the 
Hofheinz-H\"{o}velmanns-Kiltz (HHK) transform~\cite{DBLP:conf/tcc/HofheinzHK17} 
to a IND-CPA construction, which is also denoted as HQC.PKE.
Given the sound security record of HQC, much attention has been devoted to its 
side-channel cryptanalysis.
Most of the independent analyses focused mainly on power analysis or on 
extracting electromagnetic traces~\cite{wcc22,DBLP:conf/pqcrypto/GoyLG22,DBLP:conf/pqcrypto/SchambergerHRWS22,DBLP:conf/cardis/SchambergerRSW20};
nonetheless, in~\cite{DBLP:journals/tches/GuoHJLNS22} the authors described a 
timing-based key-recovery  attack against HQC, based on an implementation 
weakness in the step that used a procedure called \emph{rejection sampling} 
during randomness generation. 
As a countermeasure, a more careful implementation of the rejection sampling 
step is currently employed.\\
\indent 
To date, no timing and power- or Electromagnetic(EM)-based side-channel secure 
implementation of HQC, targeted to embedded systems, has been proposed, despite 
the real threat set by these attacks.
Moreover, we also note that most of the literature focused on timing side-channel
attacks aims at key retrieval, while it is also possible to perform a 
message-retrieval attack against the encryption primitive, in case the 
operations in the said encryption are not constant time.
Indeed, the attack mechanism mimics the common attacks on modular exponentiation, 
as the non-constant time polynomial multiplication techniques, such as the one
in the reference code by HQC, have a loop whose body contains branches depending
on the values of a secret operand.
Since, in the HQC.PKE encryption primitive, at least one of such operands is 
secret (as it is the plaintext message to be encrypted), providing constant time
multiplication algorithms is needed to secure the encryption primitive. 
Furthermore, as the HHK transform requires a execution of the HQC.PKE encryption 
function during the computation of the HQC KEM decryption, the requirement for 
constant time multiplication implementations is further strengthened.\\
\indent 
The current state of the software HQC implementations tailored to the ARM
Cortex-M$4$ microcontroller (chosen by NIST as the reference platform for
microcontroller benchmarks) includes only the one in PQM$4$~\cite{PQM4}, which is 
a public software project collecting the implementation of all NIST
fourth round candidates tailored for the Cortex M$4$ CPUs.
In the PQM$4$ project, the latest version of the HQC implementation is directly
derived from the third round reference one, without further optimizations.
As a consequence, the implementation is vulnerable to the timing side-channel 
attack by~\cite{DBLP:journals/tches/GuoHJLNS22}. Furthermore, no protection 
against power-based side-channel attacks is present.

\noindent\textbf{Our Contribution.} In this paper, we present the first fully 
constant-time and power/EM side-channel protected software implementation of HQC.KEM, tailored to Cortex-M$4$ microcontrollers. 
We provide an experimental validation of the security of our proposal via a 
battery of statistical $t$-tests, both for the EM-emissions and execution time 
side-channels.
We evaluate the timing side-channel immunity while varying the CPU frequency and 
memory latencies, to highlight the stability of our solution across different 
memory access latencies (which depend on the CPU frequency). 
Our results show that employing a FFT-based multiplication algorithm is the 
optimal choice for constant-time multiplications on Cortex-M$4$, in contrast 
with classic divide-et-impera techniques (e.g., Karatsuba) which provide a 
better tradeoff for x$86$\_$64$ computing platforms.
The new algorithmic strategy we introduce to hinder power/EM side-channel 
attacks against HQC.KEM software implementations achieves a lower than linear 
performance overhead in protecting the HQC.KEM Decapsulation primitive, 
currently a unique and welcome feature among side-channel protected 
implementations of post-quantum  KEMs.
Our approach is specifically tailored to the decapsulation of 
HQC, while applying an $n$ out of $n$ share-splitting strategy to 
Key-generation and Encapsulation primitives with a linear 
performance overhead.
